\section{A fixed-reference audit for distribution repair}
\label{sec:method}
Figure~\ref{fig:pipeline} summarizes the audit, from a fixed reference
to a controlled intervention test.
Let $P$ be the reference law and $Q_a$ the output law of configuration $a$,
which specifies the checkpoint, schedule, guidance, solver, and randomization.
Use the law $P_T$ of a resolved teacher to measure sampler-relative change
or the observed-data law $P_D$ to measure data-relative representation.
Declare the representation, regions, effect sizes, fidelity margins, comparison
family, and budget. Disjoint reference pools $R_{\rm fit}$, $R_{\rm cal}$,
and $R_{\rm eval}$ fit the regions and support score, calibrate its threshold,
and estimate reference outcomes, respectively. Candidate evaluation images
enter neither fit nor calibration. We use a diagnosis pool to choose the
intervention and a separate outcome pool to test the frozen comparison.
The audit returns a regional decision and uncertainty interval, a
conditional spread estimate, and fidelity and cost measurements for each
tested configuration. An unresolved certificate leaves open both the
presence and the absence of a practically meaningful deficit.

Let $\mathcal X$ be the image space and $\phi:\mathcal X\to\mathbb R^d$
the representation with $d$ features. Reference-fit principal component
analysis (PCA) gives a mean $m$, an orthonormal
basis $U\in\mathbb R^{d\times k}$ with up to $k=32$ retained directions,
and coordinates $z(x)=U^\top(\phi(x)-m)$. For feature regions, fit
$K=10$ $K$-means cells after scaling each coordinate by its reference-fit
standard deviation. The fixed assignment $J(x)\in\{1,\ldots,K\}$ defines
$A_j=\{x:J(x)=j\}$. Pixel, latent, DINOv2, and CLIP features offer
different views of the images~\cite{oquab2024dinov2,radford2021clip}.
For semantic regions, apply the same classifier and rejection rule to
both laws. The partition includes rejection, and rejected images stay
in the denominator of each region's probability. Held-out real
accuracy validates the labeling rule on real images. Feature fidelity
and uncurated samples help assess its interpretation on generated images.

For image $X$ under $P$, let $Y=\phi(X)$. Write $p_j=P(X\in A_j)$ and
$q_{aj}=Q_a(X\in A_j)$. Let $\mu$ be the mean of $Y$ under $P$, and let
$\mu_j$ and $\Sigma_j$ be its mean and covariance conditional on $A_j$.
Then
\begin{equation}
\operatorname{Cov}_{P}(Y)
=\sum_j p_j\Sigma_j+\sum_j p_j(\mu_j-\mu)(\mu_j-\mu)^\top.
\label{eq:totalcov}
\end{equation}
Terms with zero mass contribute zero. The first term averages within-region covariance. The second measures
separation of regional means. Either term can change with regional mass,
which explains why pooled spread alone cannot identify a deficit.

To certify a regional change, fix all events before drawing evaluation
images independently within each law. Reference and candidate draws may
be paired across laws. The guarantee assumes fixed evaluation sizes.
Pooling samples drawn under different prompts requires a sampling model
for the prompt mixture. We report prompt-level uncertainty separately.
Let $M$ count all declared region, condition, configuration, and
representation comparisons. Let $[p_j^-,p_j^+]$ and $[q_{aj}^-,q_{aj}^+]$
be exact binomial intervals with per-tail error
$\varepsilon=\alpha_{\rm occ}/(4M)$~\cite{clopper1934confidence}, where
$\alpha_{\rm occ}$ is the family error budget. The union bound yields
intervals for $q_{aj}/p_j$ in all positive-mass regions with simultaneous coverage at least
$1-\alpha_{\rm occ}$:
\begin{equation}
\left[q_{aj}^-/p_j^+,\ q_{aj}^+/p_j^-\right].
\label{eq:occupancy}
\end{equation}
For $p_j>0$, the upper endpoint is infinite when $p_j^-=0$.
Only eligible regions receive a ratio decision. An eligible region has
$p_j^-\geq\pi_{\min}$, a minimum reference mass. An upper endpoint below
the ratio threshold $\rho$ certifies \emph{depletion}; a lower endpoint
above $\rho$ certifies \emph{retention} at that resolution. Otherwise the
test is unresolved. Defaults are $\rho=0.5$ and $\pi_{\min}=0.01$.
For signed mass balance, bound $q_{aj}-p_j$ by
$[q_{aj}^--p_j^+,q_{aj}^+-p_j^-]$. The interval must lie below, above,
or within $[-\delta,\delta]$ to certify depletion, excess, or equivalence,
respectively. The digit study uses $\delta=0.03$ and charges the shared
reference bounds once (Supp.~Sec.~C.8).

With zero observations among $n_Q$ candidate draws,
$q_{aj}^+=1-\varepsilon^{1/n_Q}$. For an eligible region and
$0<\rho p_j^-<1$, certifying depletion in \cref{eq:occupancy} requires
\begin{equation}
n_Q>\frac{\log\varepsilon}{\log(1-\rho p_j^-)}.
\label{eq:planning}
\end{equation}
Rare regions and broad comparison families therefore demand more samples.
Even zero candidate counts cannot certify depletion in an ineligible region.

To measure variation separately from mass, compare the covariance traces
of the unwhitened coordinates within $A_j$:
\begin{equation}
r_{aj}=\frac{\operatorname{tr}\operatorname{Cov}_{Q_a}[z(X)\mid A_j]}
{\operatorname{tr}\operatorname{Cov}_{P}[z(X)\mid A_j]}.
\label{eq:spread}
\end{equation}
This ratio is defined when the region has positive mass under both laws
and the reference trace is positive. The trace sums coordinate variances
after centering each law separately.
A finite estimated log-ratio interval requires at least 30 samples from
each law and positive sample trace in both. Approximate influence-function intervals
wholly below 0.9 or above 1.1 flag contraction or expansion. For text-to-image,
we divide mean within-prompt trace for the candidate by that of the reference
in each representation and bootstrap whole matched prompts.

A support score $S(x)$ combines the whitened distance to the nearest
reference-fit point and the scaled PCA residual norm in quadrature.
The residual captures departures in directions discarded by projection,
and its scale is the median fit residual, floored at $10^{-6}$.
In the notation above the score is
\begin{equation}
S(x)=\sqrt{d_{\rm NN}(x)^2+
\bigl(\|\phi(x)-m-Uz(x)\|_2/s_{\rm res}\bigr)^2},
\label{eq:support}
\end{equation}
where $d_{\rm NN}(x)$ is the nearest fit-sample distance in whitened PCA
coordinates and $s_{\rm res}$ is the residual scale (Supp.~Sec.~C.1). With $n_{\rm cal}$
exchangeable reference calibration samples, its threshold $\tau$ is
the $\lceil(n_{\rm cal}+1)0.95\rceil$-th score in ascending order
\cite{vovk2013conditional}. If this rank is unavailable, the gate is unresolved.
Independent evaluations estimate $Q_a(S>\tau)-P(S>\tau)$.
Exact binomial endpoints bound this excess tail probability (Supp.~Sec.~A.3).
The gate passes if its upper bound is at most 0.05 and fails if its lower
bound exceeds 0.05. Otherwise it is unresolved. The feature-only default
gives occupancy and support error budgets of 0.025 each. The approximate
spread intervals are outside that exact error guarantee. Signed-balance
and routing follow-ups allocate 0.025/0.0125/0.0125 to mass/fidelity/support.
CIFAR selection also checks teacher precision, text-to-image checks alignment
and learned preference, and digits use classifier consensus and feature
precision. Expansion with a failed support gate is off-reference expansion.

For the same semantic event, let $p_{D,j}=P_D(A_j)$ and $p_{T,j}=P_T(A_j)$.
When both are positive, $q_{aj}/p_{D,j}=(q_{aj}/p_{T,j})(p_{T,j}/p_{D,j})$
separates sampler--teacher from teacher--data mass differences.
Teacher-relative deficits motivate changing the starting noise with the
denoiser, grid, and call budget fixed, or changing the grid or solver under
the same conditioning. A validated data-relative imbalance that persists
across qualified samplers motivates comparing ordinary and targeted
continuations from the same parent at equal update budgets. Persistence across adequate
samplers motivates a training comparison because the tested
longer samplers did not resolve the deficit. Freeze the
comparison after diagnosis and evaluate its targeted outcome, fidelity,
coverage, and cost on independent samples. Final-sample deficits can have
several causes. These controls test the effect of one specified change.
